\subsection{映入赋范空间的映射空间}

更特别地，考虑 Y 是非离散赋值除环 K 上的赋范向量空间的情形（第 IX 章，§ 3，第 3 号）。用 \(||y||\) 表示 \(y \in Y\) 的范数。这时集合 \(\mathcal{F}(X;Y) = Y^{\mathbf{X}}\) 典范地赋有一个 K 向量空间结构。映射 \(u: X \to Y\) 有界当且仅当实值函数 \(x \to ||u(x)||\) 在 X 中有界。若 \(u, v\) 是 X 到 Y 中的有界映射，则显然 \(u + v\) 与 \(\lambda u (\lambda \in K)\) 都有界；换句话说，\(\mathcal{B}(X;Y)\) 是 \(\mathcal{F}(X;Y)\) 的一个向量子空间。此外，\(||u|| = \sup_{x \in X} ||u(x)||\) 是 \(\mathcal{B}(X;Y)\) 上的一个范数；因为它满足三角不等式，且 \(||u|| = 0\) 蕴涵 \(u = 0\) ，又对每个 \(\lambda \in K\) 我们有

\[
| | \lambda u | | = \sup _ {x \in \mathbf {X}} | | \lambda u (x) | | = \sup _ {x \in \mathbf {X}} | \lambda |. | | u (x) | | = | \lambda |. \sup _ {x \in \mathbf {X}} | | u (x) | | = | \lambda |. | | u | |.
\]

此外，立即可以验证，由这个范数在 \(\mathcal{B}(\mathbf{X};\mathbf{Y})\) 上定义的一致结构就是一致收敛的一致结构。除非明确说明相反的情形，每当把 \(\mathcal{B}(\mathbf{X};\mathbf{Y})\) 看作赋范空间时，所指的都是上面定义的范数。

命题 3. 若赋范空间 Y 完备，则每个级数 \((u_n)\) （由 X 到 Y 中的有界映射组成，在赋范空间 \(\mathcal{B}\) (X; Y) 中绝对收敛，即满足 \(\sum_{n=0}^{\infty}||u_n|| < +\infty\) ；参见第 IX 章，§ 3，第 6 号）都在 X 中一致收敛。

因为 \(\mathcal{B}(\mathbf{X};\mathbf{Y})\) 完备（第 1 号，命题 2 的推论 1），结论由第 IX 章，§ 3，第 6 号，命题 11 以及一致收敛级数的定义得出。

\[
\text { Remark. } \quad 1) \text {   If   } \sum_ {n = 0} ^ {\infty} \| u _ {n} \| <   + \infty , \text {   then   } \sum_ {n = 0} ^ {\infty} \| u _ {n} (x) \| \leqslant \sum_ {n = 0} ^ {\infty} \| u _ {n} \| <   + \infty
\]

对每个 \(x \in \mathbf{X}\) 成立；换句话说，对每个 \(x \in \mathbf{X}\) ，以 \(u_n(x)\) 为通项的级数在空间 \(\mathbf{Y}\) 中绝对收敛。逆命题不成立。为避免一切混淆，我们有时说以 \(u_n\) 为通项的级数是正规收敛的，意思是以 \(\| u_n \|\) 为通项的级数收敛。一个级数可以在 \(\mathbf{X}\) 中一致收敛而非正规收敛；例如级数 \((u_n)\) （在空间 \(\mathcal{B}(\mathbf{R}, \mathbf{R})\) 中，定义如下：\(u_n(x) = (1/n) \sin x\) 若 \(x \in [n\pi, (n + 1)\pi]\) ，否则 \(u_n(x) = 0\) ）就是这种情形。

当 Y 是非离散赋值域 K 上的赋范代数（第 IX 章，§ 3，第 7 号）时，\(\mathcal{B}(\mathbf{X};\mathbf{Y})\) 是一个 K 代数，且范数 \(||u||\) 与代数结构相容，因为

\[
\begin{array}{r l} | | u v | | = \sup _ {x \in \mathbf {X}} | | u (x) v (x) | | & \leqslant \sup _ {x \in \mathbf {X}} | | u (x) | |. | | v (x) | | \\ & \leqslant \sup _ {x \in \mathbf {X}} | | u (x) | |. \sup _ {x \in \mathbf {X}} | | v (x) | | = | | u | |. | | v | |. \end{array}
\]

这样，\(\mathcal{B}(\mathbf{X};\mathbf{Y})\) 现在是 \(\mathbf{K}\) 上的一个赋范代数。

命题 4. 设 \(\mathbf{X}_i\) （ \(1 \leqslant i \leqslant n\) ）与 \(\mathbf{Y}\) 是非离散赋值除环 \(\mathbf{K}\) 上的赋范向量空间，并设 \(\mathbf{X} = \prod_{i=1}^{n} \mathbf{X}_i\) 。则 \(\mathbf{X}\) 到 \(\mathbf{Y}\) 中的一切多重线性映射所成之集在空间 \(\mathcal{F}_s(\mathbf{X}; \mathbf{Y})\) 中是闭的。

这个集由满足所有下列关系的 \(u \in \mathbf{F}(\mathbf{X};\mathbf{Y})\) 组成

\[
\begin{array}{r l} u (x _ {1}, \dots , x _ {i} ^ {\prime} + x _ {i} ^ {\prime \prime}, \dots , x _ {n}) & = u (x _ {1}, \dots , x _ {i} ^ {\prime}, \dots , x _ {n}) \\ & \quad + u (x _ {1}, \dots , x _ {i} ^ {\prime \prime}, \dots , x _ {n}), \\ u (x _ {1}, \dots , \lambda x _ {i}, \dots , x _ {n}) & = \lambda u (x _ {1}, \dots , x _ {i}, \dots , x _ {n}) \end{array}\tag{1}
\]

（ \(1 \leqslant i \leqslant n, x_i, x_i', x_i''\) 为 \(\mathbf{X}_i\) 的任意元素，\(\lambda\) 为 \(\mathbf{K}\) 的任意元素）；由于关系（1）的两端都是 \(u\) 在 \(\mathcal{F}_s(\mathbf{X}; \mathbf{Y})\) 上的连续函数（ \(\S\) 1，第 2 号，注 6），结论成立（第 I 章，§ 8，第 1 号，命题 2）。

命题 5. 在命题 4 的假设下，集 \(\mathcal{L}(\mathbf{X}_1,\ldots ,\mathbf{X}_n;\mathbf{Y})\) （由 \(\mathbf{X}\) 到 \(\mathbf{Y}\) 中的连续多重线性映射组成）在 \(\mathcal{F}(\mathbf{X};\mathbf{Y})\) 中关于有界收敛拓扑是闭的；若 \(\mathbf{Y}\) 完备，则它关于有界收敛的一致结构是完备的。

因为若 \(\mathfrak{S}\) 是 \(\mathbf{X},\mathcal{L}(\mathbf{X}_1,\ldots ,\mathbf{X}_n;\mathbf{Y})\) 就是 \(\mathbf{X}\) 到 \(\mathbf{Y}\) 中一切多重线性映射所成之集与集 \(\mathcal{B}_{\mathfrak{S}}(\mathbf{X};\mathbf{Y})\) 的交（第 IX 章，§ 3，第 5 号，定理 1）；于是结论由命题 4 与命题 2 的推论 1 得出。

在本小节的其余部分，K 表示一个非离散赋值域。

于是 \(\mathcal{L}(\mathbf{X}_1, \ldots, \mathbf{X}_n; \mathbf{Y})\) 是 \(\mathcal{F}(\mathbf{X}; \mathbf{Y})\) 的一个向量子空间。设 \(\mathbf{B}\) 是 \(\mathbf{X}\) 中的单位球，即所有 \((\boldsymbol{x}_i)_{1 \leqslant i \leqslant n}\) （满足 \(\sup_{1 \leqslant i \leqslant n} ||\boldsymbol{x}_i|| \leqslant 1\) ）所成之集。则映射 \(u \to u|\mathbf{B}\) （由 \(\mathcal{L}(\mathbf{X}_1, \ldots, \mathbf{X}_n; \mathbf{Y})\) 到 \(\mathcal{B}(\mathbf{B}; \mathbf{Y})\) 中）是单射；此外，\(\mathcal{B}(\mathbf{B}; \mathbf{Y})\) 上的一致收敛的一致结构在这个映射下的原像是 \(\mathcal{L}(\mathbf{X}_1, \ldots, \mathbf{X}_n; \mathbf{Y})\) 上的有界收敛的一致结构。因为 \(\mathbf{X}\) 的每个有界子集都含于形如 \(\mu\mathbf{B}\) 的集（对某个 \(\mu \in \mathbf{K}^*\) ），且若 \(u\) 是 \(\mathcal{L}(\mathbf{X}_1, \ldots, \mathbf{X}_n; \mathbf{Y})\) 的一个元素，则说 \(||u(z)|| \leqslant a\) 对一切 \(z \in \mu\mathbf{B}\) 成立等价于说 \(||u(z)|| \leqslant a/|\mu|^n\) 对一切 \(z \in \mathbf{B}\) 成立。容易验证，数

\[
\left| \left| u \right| \right| = \sup _ {z \neq 0} \frac {\left| \left| u (z) \right| \right|}{\left| \left| z \right| \right|}
\]

是 \(\mathcal{L}(\mathbf{X}_1, \ldots, \mathbf{X}_n; \mathbf{Y})\) 上的一个范数，并定义这个集上的有界收敛的一致结构，而且显然我们有

\[
\left| \left| u \left(x _ {1}, \dots , x _ {n}\right) \right| \right| \leqslant \left| \left| u \right| \right|. \left| \left| x _ {1} \right| \right| \dots \left| \left| x _ {n} \right| \right|.\tag{2}
\]

除非明确说明相反的情形，每当把 \(\mathcal{L}(\mathbf{X}_1,\ldots ,\mathbf{X}_n;\mathbf{Y})\) 看作赋范空间时，所指的都是上面定义的范数。

命题 6. 多重线性映射

\[
(u, x _ {1}, \dots , x _ {n}) \rightarrow u (x _ {1}, \dots , x _ {n})
\]

（由赋范空间 \(\mathcal{L}(\mathbf{X}_{1},\ldots,\mathbf{X}_{n};\mathbf{Y})\times\mathbf{X}_{1}\times\cdots\times\mathbf{X}_{n}\) 到 Y 中）是连续的。

这是不等式（2）的直接推论（第 IX 章，§ 3，第 5 号，定理 1）。

命题 7. 设 X、Y、Z 是 K 上的三个赋范空间。由赋范空间 \(\mathscr{L}(X, Y; Z)\) 到（X 到 \(\mathscr{L}(Y; Z)\) 中的线性映射所成的）空间的典型映射，它把每个 \(u \in \mathscr{L}(X, Y; Z)\) 映为映射 \(x \to u(x, .)\) ，是 \(\mathscr{L}(X; Y; Z)\) 到 \(\mathscr{L}(X; \mathscr{L}(Y; Z))\) 上的一个等距。

这由定义以及关系式

\[
\sup _ {| | \boldsymbol {x} | | \leqslant 1} \left(\sup _ {| | \boldsymbol {y} | | \leqslant 1} | | u (\boldsymbol {x}, \boldsymbol {y}) | |\right) = \sup _ {| | \boldsymbol {x} | | \leqslant 1, | | \boldsymbol {y} | | \leqslant 1} | | u (\boldsymbol {x}, \boldsymbol {y}) | |.
\]

命题 8. 设 X、Y、Z 是 K 上的三个赋范空间。双线性映射 \((u, v) \to v \circ u\) （由 \(\mathcal{L}(X; Y) \times \mathcal{L}(Y; Z)\) 到 \(\mathcal{L}(X; Z)\) 中）是连续的。

因为若 \(u \in \mathcal{L}(\mathbf{X};\mathbf{Y})\) 且 \(v \in \mathcal{L}(\mathbf{Y};\mathbf{Z})\) ，我们有

\[
\left| \left| v \circ u \right| \right| \leqslant \left| \left| u \right| \right| \cdot \left| \left| v \right| \right|,\tag{3}
\]

因为对一切 \(\pmb{x} \in \mathbf{X}\) ，由（2）我们有 \(||v(u(\pmb{x}))|| \leqslant ||v|| \cdot ||u(\pmb{x})|| \leqslant ||v|| \cdot ||u|| \cdot ||\pmb{x}||\) 。

特别地，在集 \(\mathcal{L}(\mathbf{X})\) ——赋范空间 \(\mathbf{X}\) （在 \(\mathbf{K}\) 上）的连续自同态所成之集——上，范数 \(||u||\) 与 \(\mathcal{L}(\mathbf{X})\) 的 K 代数结构相容。

注。2) 集 \(\mathcal{L}(\mathbf{R}^{m};\mathbf{R}^{n})\) （由 \(\mathbf{R}^{m}\) 到 \(\mathbf{R}^{n}\) 中的线性（必为连续的）映射组成）可等同于矩阵所成之集 \(\mathbf{M}_{n,m}(\mathbf{R})\) （\(n\) 行 \(m\) 列，系数在 \(\mathbf{R}\) 中），从而可等同于 \(\mathbf{R}^{mn}\) ；在 \(\mathcal{L}(\mathbf{R}^{m};\mathbf{R}^{n})\) 上，有界收敛（关于 \(\mathbf{R}^{m}\) 上的欧氏度量）、紧收敛以及逐点收敛的一致结构这时都等同于 \(\mathbf{R}^{mn}\) 上的加法一致结构。取 \(x = (x_{i})\in \mathbb{R}^{n}\) 的范数为

\[
\| \boldsymbol {x} \| = \sup _ {i} | x _ {i} |,
\]

并设 \((\pmb{e}_j)\) 是 \(\mathbf{R}^m\) 的标准基；若 \(u\) 与 \(v\) 是 \(\mathbf{R}^m\) 到 \(\mathbf{R}^n\) 中的两个线性映射，满足 \(\| u(\pmb{e}_j) - v(\pmb{e}_j)\| \leqslant \varepsilon\) 对 \(1 \leqslant j \leqslant m\) ，则我们有 \(|\alpha_{ij} - \beta_{ij}| \leqslant \varepsilon\) 对每对 \((i,j)\) \([U = (\alpha_{ij})\) 与 \(V = (\beta_{ij})\) 分别是 \(u, v\) 的矩阵{]}；反之，若这些不等式成立，则我们有 \(\| u(\pmb{x}) - v(\pmb{x})\| \leqslant ma\varepsilon\) 对每一点 \(\pmb{x}\) （一个以 0 为中心、以 \(a\) 为边长、位于 \(\mathbf{R}^m\) 中的方体的点）成立。

